
d'une
structure de @@M0@@-anneau, les relations (2.4.2) passant au quotient.

Toute intersection de @@M1@@-idéaux est un @@M2@@-idéal. En particulier, il existe un plus petit @@M3@@-idéal contenant une famille donnée d'éléments @@M4@@~; d'après les relations (2.4.2), c'est l'idéal (ordinaire) engendré par les @@M5@@, @@M6@@.

Soient toujours @@M7@@ un @@M8@@-anneau, @@M9@@ une famille d'indéterminées, @@M10@@, et soit @@M11@@ une famille d'éléments de @@M12@@ : soit @@M13@@ le @@M14@@-idéal engendré par les @@M15@@.

\textbf{Définition 4.6.} \uline{Le} @@M16@@-\uline{anneau} @@M17@@ \uline{sera appelé} @@M18@@-anneau engendré sur @@M19@@ par la famille de générateurs @@M20@@ soumis aux relations @@M21@@.

On peut également se donner des relations entre les @@M22@@, @@M23@@ et faire la même construction, en remarquant qu'un idéal est un @@M24@@-idéal si et seulement si il est stable pour les opérations @@M25@@, @@M26@@, comme il résulte des relations 3.2.

Soit @@M27@@ une @@M28@@-@@M29@@-algèbre, @@M30@@ une famille d'éléments de @@M31@@. Les @@M32@@ étant des polynômes par rapport aux indéterminées @@M33@@, @@M34@@, on peut substituer à ces indéterminées les @@M35@@, @@M36@@~; si l'élément obtenu est nul, on dira que les @@M37@@ satisfont les relations @@M38@@. La @@M39@@-@@M40@@-algèbre @@M41@@ est alors caractérisée par la propriété universelle suivante :

\textbf{Proposition 4.7.} \uline{Pour toute} @@M42@@-@@M43@@-\uline{algèbre} @@M44@@ \uline{et toute famille} @@M45@@ \uline{d'éléments de} @@M46@@, \uline{satisfaisant les relations} @@M47@@, \uline{il existe un} @@M48@@-@@M49@@-\uline{homomorphisme et un seul} @@M50@@ \uline{tel que pour tout} @@M51@@, on ait @@M52@@.

En effet, le noyau de l'homomorphisme @@M53@@ défini en 4.2 contient les @@M54@@, donc @@M55@@ puisque c'est un @@M56@@-idéal~; @@M57@@ passe donc au quotient et définit 
